\section{Inherited flight channels and derived expressions}
The ordinary result schema covers time and position; altitude and distance; velocity and acceleration; Mach number and Reynolds number; orientation, angle of attack and angular rates; mass, CG and inertias; aerodynamic coefficients and CP; thrust and drag; atmospheric pressure, temperature and density; wind and gravity; and stability and event-related quantities. Not every branch, stepper or imported dataset supplies every channel. Channel availability is a property of the data and producer, not a promise made by a plot menu.

Summaries such as maximum altitude, maximum speed, deployment velocity and flight time are extracted from the recorded flight and its event context. Their definitions matter when a run aborts, a branch separates, or a time limit ends the run before impact. The GUI's simulation table is a compact view of these values; its table export is distinct from exporting all time-series samples.

The expression subsystem uses the bundled \code{de.congrace.exp4j} parser and OpenRocket's own function extensions. It resolves flight-channel symbols and supports subexpressions that select a value at a time or a range of samples. The range syntax enables functions such as mean, minimum, maximum, variance, standard deviation, RMS, trapezoidal integration, and nearest-value time. Ordinary arithmetic, powers and trigonometric functions are also available. The expression editor provides names, symbols, unit labels, validation and a function list.

For example, a derived expression can combine recorded velocity components or summarize a chosen time interval. This operates on the data exposed to the expression evaluator; it does not define a new equation of motion. Unit labels for custom expressions do not constitute a full symbolic dimensional-analysis system. The author must keep the input symbols, internal units and chosen output unit consistent. A formula referring to a missing channel cannot create the missing measurement.

Range operations deserve special care because the integrator's accepted time intervals can vary. The expression subrange machinery supplies its sampled array and associated interval metadata to functions such as \code{trapz}; it is that representation that is integrated. For precision-sensitive independent analysis, export the actual time column together with the needed channels and state the interpolation/integration procedure.

\sourceline{\source{core/src/main/java/info/openrocket/core/simulation/customexpression/CustomExpression.java}{expression parsing and evaluation}; \source{core/src/main/java/info/openrocket/core/simulation/customexpression/Functions.java}{extra functions}; \source{swing/src/main/java/info/openrocket/swing/file/SimulationTableCSVExport.java}{simulation summary export}.}

\section{CSV controls and component-analysis output}
\class{CSVExport} writes selected flight-branch data or component-analysis sweep data. Export settings include the chosen quantities and units, separator, precision, exponential notation, comments and header/event information. Unit conversions occur at export; the recorded branch remains in its internal representation. For a reproducible table, retain the time or sweep-domain column, quantity names, units and enough context to identify the branch and configuration.

The component-analysis overload exports quantities associated with both a component and a sweep coordinate. Its independent variable might be Mach or angle of attack rather than time. A series of CP values in this export is therefore a controlled parameter sweep, not automatically the CP history of the simulated flight. Likewise, per-instance and installed-total drag should be labeled explicitly when comparing components with different repetition counts.

\sourceline{\source{core/src/main/java/info/openrocket/core/file/CSVExport.java}{CSVExport}; \source{core/src/main/java/info/openrocket/core/componentanalysis/CADataType.java}{component-analysis output types}.}

\section{Fabrication templates and design reports}
The print-selection model exposes seven inherited output families: design report, parts detail, nose-cone templates, transition templates, centering-ring templates, fin templates, and fin-marking guides. Some choices are stage-specific. \class{PrintController} dispatches to printable shapes and visitor strategies, with paper size, orientation, margins and layout handled by the printing layer.

Fin templates turn the modeled outline into a physical-scale pattern; centering-ring templates include the relevant ring and tube geometry; transition and nose outputs derive patterns or profiles appropriate to those component implementations. The fin-marking guide helps lay out angular fin positions on the body. These outputs are construction aids derived from the supported component geometry, not a general CAD machining or unfolded-composite manufacturing system.

The design report combines illustrations and design information with applicable simulation information. Its simulation worker is a separate part of report preparation, so generating a report can require computations beyond drawing the current screen. Printed values must be interpreted with their configuration and units, and fabrication templates must retain their intended scale through the PDF viewer/printer settings. The geometry generator's true-size output can still be rescaled by an external print dialog.

\sourceline{\source{swing/src/main/java/info/openrocket/swing/gui/print/OpenRocketPrintable.java}{printable families}; \source{swing/src/main/java/info/openrocket/swing/gui/print/DesignReport.java}{DesignReport}; \source{swing/src/main/java/info/openrocket/swing/gui/print/FinMarkingGuide.java}{FinMarkingGuide}; \source{swing/src/main/java/info/openrocket/swing/gui/print/PrintSimulationWorker.java}{report simulation worker}.}
